\documentclass[12pt]{article}  
\newcommand{\say}[1]{``#1''}   
\newcommand{\nsay}[1]{`#1'}   
\usepackage{endnotes}   
\newcommand{\B}{\backslash{}}   
\renewcommand{\,}{\textsuperscript{,}}   
\usepackage{setspace}   
\usepackage{tipa}   
\usepackage{hyperref}   
\renewcommand{\title}[2]{\section{\href{#1.html}{#2}}}  
\newcommand{\published}[1]{First Published: #1}  
\newcommand{\draft}[1]{\section{Draft #1}}  
\newcommand{\dailyreflect}[1]{\section{Daily Reflection: #1}}  
\begin{document}  
\title{revamping-life}{On Revamping My Life}  
\published{2026 August 17}

\draft{2: 17 August 2026}

Woot! I got trapped in the \say{let's go through the 4tW\footnote{for those who don't know what this acronym is, I'm sure I've written about it before. It's a site that motivates writing, at least for me} files and sort them} because I realized most of the files I had tagged as \say{blog} were, in fact, posted blogs.  
So, I didn't end up publishing it yesterday, and I didn't end up making my daily reflection or anything else in the list from yesterday.   
For my own sake, the to do list is:

\begin{itemize}

\item Figure out when to blog, both for now and when I get a new schedule again  
\item Revise the template so that it contains my ongoing posts instead of upcoming posts\footnote{because I haven't been struggling with idea generation lately, probably because I'm still not in a heavy routine}  
\item Revise the daily reflection for better sets of questions.

Still need to think about that

\item Figure out my home.

Did make progress here: cleared off my large table, did some vague food prep\footnote{pre-cut pork chop/loin (I forget which and honestly think it might have been both) were on sale at Costco this week, so I bought the smallest possible packages (because it's a flat \$x off the package, so the percentage discount is best for the lightest) and fried two two days ago (mmmm two two and grammar works) and then baked the rest lasterday (which is different than yesterday). Recipe: put pork on parchment paper on metal tray (in my case I think it's a pizza steel), turn oven to 400, generously salt both sides of each pork, when oven preheated, cook for approx 20 minutes, then turn broiler on high for approx 4 minutes, then turn oven off and let rest in there for a few minutes, then pull out and let cool before bagging (or just nomming)}\,\footnote{oh, I also chopped the two pounds of celery to freeze (because I love being able to just grab some pre-chopped onion or celery (or sometimes carrot) when I want to cook)}, and I also assembled the couch I was given.

The couch is super nice! It's a sectional that comes into three pieces\footnote{which makes transport miles away easier!}: the two seats (both recliners!) and a center console which is two cupholders and what I assume is meant as a trash or ice box.  
It's a lil curved which is fun, and it holds the projector at a decent height, especially if I have a book to stabilize it underneath. I still do want to hang the projector somehow, but need to figure that out before I can.

\item Career aspirations: what are they, how do I progress in them?

\item Web serials: What's stopping me, do I actually want them?

\item Music: how do I make myself do them?

I did do the helpful thing last night of restringing my guitar, but I also have the issue now that I lost my wire cutters, so the strings are long and unwieldy again.

\item Finish RebelFit

\item Stretching into a daily routine

\end{itemize}

And then we can add the new ones from today:

\begin{itemize}

\item Clean my digital footprint.

 As the upcoming posts section for the past many posts implied, my inbox is now just under 152,000 emails large, and that includes me deleting a bunch recently and having good spam and automatic filters.  
The oldest email I see is from 2012, which is strange because it's definitely not on my oldest email account.  
Aha, the oldest email account has another few months of emails, just not in the inbox\footnote{and would add approximately another 40000 emails to the inbox if that list was included}.  
I'm sure there are older messages somewhere, because Gmail was released in 2004 and the oldest emails are still only from 2008, but that's quite possibly the oldest emails that still exist.

Anyways, I don't like that I have infinite useless messages, because it makes the messages harder to search and also like it's a waste of space and time.  
I kept them for some reason that I'd still like to blog about\footnote{and it's likely to become an upcoming blog post}

\item Figure out what other forms of writing I want to be doing

I'm not sure if the web serial is in this phase of my life or not, but if not, then I need to have some other form of writing.

\item Get back into touch type practice.

At this point, I'm fairly convinced that, while QWERTY may not be the absolute ideal keyboard layout, short of learning stenography, it's likely to be my best bet, because the opportunity cost of learning a different layout is much too high and it doesn't seem like there's good data for people being significantly faster on average.  
There's some strain stuff, but since I never want to get to a point that I'm typing to physical failure, that's not likely to be an issue for me for years, at a bare minimum.

Unrelated: \say{anyways} is apparently the older attested form, according to Merriam Webster. Its first appearance is \href{https://www.merriam-webster.com/dictionary/anyways#word-history}{thirteenth century}, but anyway only dates to \href{fourteenth}{https://www.merriam-webster.com/dictionary/anyway#word-history}.  
MW also agrees that it's a real, though dialectical\footnote{presumably using their second definition of the word} \href{word}{https://www.merriam-webster.com/grammar/lets-talk-about-anyways}.  
That's fun, anyways.\footnote{and now adding the word to my dictionary (which means right clicking the red highlighted word (triple clicking, since Mac), and clicking \say{Add to Dictionary}}.

\item Get a regular physical writing practice established.

I don't, at this point, really care if the practice is poetry, journaling, notes, or really anything, so long as I am writing again.  
I miss having a distinctive and pretty\footnote{at least in my own view} hand, so I'd also like to work on bringing one of those again.  
I think that I prefer print, if only because cursive is something that children learn much later, and it's also something that I feel I waste a lot of time trying to reread.

\item Read through the old writing I have produced.

I printed and bound a full copy of my web serial at some point, and even though I don't remember where it went\footnote{I have a memory of giving it to either my partner or a close friend}, I would like to read through that.

\item Get into a regular embroidery habit.

Even if I'm not doing the huge project, which is still in the design stage, it would be nice to be able to make small embroidered objects to give to people I love and care about.  
Since the free-hand aspect of embroidery is what appeals to most people, that's probably something that I'd also like to get into.  
Something to consider, especially because random fabric is miles away cheaper than nice even-weave, especially Aida and its off-brand clones.

Huh, never mind.

Aida as a fabric at least dates back to Verdi's opera of the same name, so is from the late nineteenth century.  
I always assumed a brand name or trademark.  
Matty fabric may still be cheaper, though.

\item Get into some sort of prayer habit.

I have found myself in a spiritually empty place for a hot second now.

That is not a feeling I like or want to keep.

Best bet I can think of to get there is to pick up something Catholic and something Jewish.  
I have a history of Catholic prayer, so it might be good for me to pray the Jesus prayer a few times\footnote{Lord Jesus Christ, Son of G-d, have mercy on me, a sinner}, or perhaps something less inherently self-denigrating, given the current mental battles.  
At the very least I might want to get back into praying before meals.

The prayer for St. Michael might be a good one?  
Let's see, the text goes something like:\footnote{from memory, given that I did set it to music and win an award for the setting, I should remember it (to say nothing of the many times I said it myself)}

St. Michael the Archangel\footnote{I think capital}, defend us in battle. Be our protection against the wickedness and snares of the devil. May G-d rebuke him, we humbly pray, and do thou, oh prince of the heavenly host, by the power of G-d, cast\footnote{or thrust, or many other words} into Hell Satan and all the evil spirits\footnote{or ghosts, which is the version I used in my setting because syllables} that prowl about the world, seeking the ruin of souls.

Looking at \href{https://www.ewtn.com/catholicism/devotions/prayer-to-st-michael-the-archangel-371}{EWTN}, they use \say{defense} instead of \say{protection}, \say{heavenly hosts} plural not singular, \say{thrust} instead of \say{cast} (which I also have used in the past at times).  
Woot! I do have a memory still.

Apparently there are three daily Jewish prayers, though that seems as though it may be more of a communal activity.  
Alas, I have not explored my Jewish side as much as I might.

\item Jewish identity: explore.

There's a weekly Shabbat service from a synagogue\footnote{from Greek via Latin to Old French to Middle English for gathering} that I have a friend who cantors for.\footnote{huh wild that only when speaking about Jewish things do I find myself hearing a pretentious Boston voice telling me to never end a sentence with a preposition (because of the old joke \say{guy asks \nsay{where's the X at}, reply: \nsay{my good sir, don't you know that it is improper English for to end a sentence with a preposition? As a Harvard professor, I feel it is my duty to inform you of this fact.} first: \nsay{fine. Where's the X at, (insert preferred slur (people don't like that I use slur to mean any term meant to degrade a person, which is potentially fair)/ insulting word)?}}, if I had to guess.). That's fun to realize I still think of my father's jokes from when I was small}  
I could try to start attending that, if nothing else, because while I do find that I love a rigid belief system, it's not what I actually find solace in.

\end{itemize}

Anyways, that's a large enough list of items for now! I guess this could work well for something that I have as a daily reflection:

\begin{enumerate}

\item What time are you starting this reflection? How's this time feel?

\item What time are you ending this reflection? How's the time feel? How'd reflecting go?

\item Blogging: when did you last blog? How did the time feel? How long did you spend?

\item Blogging: Any changes needed to the template? If so, why haven't you made them?

\item Daily reflection: Any questions needed to this template? If so, what and why?

\item Home: What's one small thing you can do today to improve your home setup?

\item Home: What's one medium thing you can do today to improve the home setup?

\item Home: What's one large thing that home setup needs right now? Is there a way to break it into small and medium tasks you can do?

\item Career aspirations: Any changes?

\item (assuming I want to do web serial(s)):

\begin{itemize}

\item New serial(s): what's one large thing that the new serial(s) needs right now? Can you break it to small/medium tasks?

\item New serial(s): What's at least one medium thing that the new serial(s) need)s(\footnote{inverse parentheses!. That's dumb, never use again}? Can it or they be broken into smaller tasks?

\item New serial(s): What's at least one small thing that the new serial(s) needs? Is that something you can do right now? If not, what's something small you can do right now?

\item Old serial(s): What's the status of plotting out the rest of the book?

\item Old serial(s): What's the status of writing the rest of it?

\end{itemize}

\item Music: Have you sung a scale today? If not, do a vocal warmup (hmm change to just \say{have sung?} in the next one)

\item Music: Guitar? If not, do right and left hand flexibility exercise now, at least for a minute or so.

\item Music: Composition? What is the smallest possible piece of composition you could do right now that would feel like it \say{counts} to you? Do that now.

\item RebelFit: What's in the way of it being finished? Can you break it into large, medium, small tasks?

\item RebelFit: What's at least one large task? Can you break it into medium and small tasks?

\item RebelFit What's at least one medium task? Can you break it or them into small tasks?

\item RebelFit: What's at least one small task? Can it be done within two minutes? If not, it's not small. If so, go do it.

\item Stretching: Go stand and touch your toes for three deep breaths.

\item Stretching: Do you have a stretching routine for 10 minutes? If not, what is missing from it? If so, when's the last time you did it?

\item Stretching: same for 15 and/or 30 minutes(?)

\item Stretching: Same for 5 minutes.\footnote{apple watch only counts 5 minutes and up as one workout, so that's the smallest unit of time I'll ever have}

\item Cleaning inbox: What's the current mail count? Go move or delete 25 emails, based on whether they need to be kept or not.

\item Blog: What's the biggest block to finishing one of your \verbatim{\href{blogs-to-write}{in-progress}} (when I make that post, it's going to be its own post, I think? That makes sense to me, and then I can just reference the live page for it?) blogs?

\item Touch type: Have you spent five minutes touch typing today? (n.b. for the reader, I really love \href{keybr.com}{keybr} for at least the 26 normal letters, haven't tried it much with punctuation or numbers or non-English uses\footnote{e.g. programming, score writing})

\item What other forms of writing are you doing right now? Do them.

\item Physical Journal Writing (if that's what I choose): When did you last do this? How did the time feel? If it was more than 24 hours ago, why hasn't it happened again?

\item Physical Poem Writing (if that's what I choose): When did you last write a poem? How did the time feel? If more than 24 hours ago, why hasn't happened again? What poem type was it, how did that feel and go? What aids and tools did you need?

\item Physical Writing: How's the new hand looking? Your current goals were that it is: printed, pretty (other goals), any progress on that front? Think using the n-gram mapper you made would help?

\item Blogging: Have you written and posted your \say{blogs-im-writing} blog post yet? Why not?

\item Blogging: Have you written and posted your \say{do a single time tasks} blog post yet? Why not? (this includes things like \say{make an n-graph for 1-5 (or more maybe) characters so that when writing you know relative frequencies to hit for optimal growth, assuming all things need equal work}, and \say{make the blog meta-data'd\footnote{that cannot be the correct spelling}} I'm now adding the 30, 15, 10, 5 minute stretch patterns to it.)

\item Reading old blog posts: (I'm going to plan to start from the beginning and add notes to a new blog post) did you read your next blog post and (depending if I can add meta-data or not) add relevant tags? Did you add a summary of it to the notes post? If not, do that.

\item Reading old journals: What journal are you on? What are you reading through now? Is it worth transfer to another medium (e.g. a good poem, a song you wrote, something to remember)?

\item Reading other old writing: What writing are you reading? Anything worth maintaining? If not, why didn't you throw it away?

\item Reading old web serial: What chapter are you up to? Did you add to the (not being posted here because I want at least the thinnest possible thread of separation between that identity and me, at least the me that posts here) notes document?

\item Regular embroidery practice: What are you doing? What's five minutes of embroidery going to have looked like? (rewrite) Was the estimate right? (Maybe also add a question for counted and non-counted thread)

\item Prayer habit: When's the last time you spent in prayer? How'd the timing feel? How'd the length feel? If more than 8 hours ago, why haven't done since?

\item Jewish Identity: How's that going?

\end{enumerate}

Realizing that I want a different order during the Home set, I'll copy to below and make the final list (for transparency, I then rewrote the list as it existed. All items after home on the above list were added after that. As a result, I started leaving notes to myself so that I'd remember what to do in the newer version).

\begin{enumerate}

\item What time are you starting this reflection? How's this time feel?

\item Blogging: when did you last blog? How did the time feel? How long did you spend?

\item Blogging: Any changes needed to the template? If so, why haven't you made them?

\item Daily reflection: Any questions needed to this template? If so, what and why?

\item Home: What's one large thing that home setup needs right now? Is there a way to break it into small and medium tasks you can do?

\item Home: What's at least one medium thing you could do today to improve your home setup? Is there a way to break it or them into small tasks?

\item Home: What's at least one small thing you can do today to improve the shape of your home? If multiple, what is the one you're going to pause right now to do, and what is the most important of the remaining?

\item Career aspirations: Any changes?

\item If doing serials, look above for questions

\item Music: Have you sung today? If not, do a vocal warmup

\item Music: Have you played guitar today? If not, do at least one hand's flexibility exercise.

\item Music: What is the smallest possible piece of composition you could do right now that would feel like it \say{counts} to you? Do that now if you haven't already composed today.

\item RebelFit: What's in the way of it being finished? Can you break it into large, medium, small tasks?

\item RebelFit: What's at least one large task? Can you break it into medium and small tasks?

\item RebelFit What's at least one medium task? Can you break it or them into small tasks?

\item RebelFit: What's at least one small task? Can it be done within two minutes? If not, it's not small. If so, go do it.

\item Stretching: When were your last three stretches? What routine (lengths) were they? How did the time feel for stretching? How were the routines? If it's been more than 36 hours since three have been done, what's in the way?

\item Cleaning inbox: What's the inbox count? Go move or delete 25 emails, based on whether they need to be kept or not. New inbox count?

\item Blogging: Have you written and posted your \say{blogs-im-writing} blog post yet? Why not?

\item Blogging: Have you written and posted your \say{do a single time tasks} blog post yet? Why not?

\item Blogging: Have you written and posted your \say{this is literally just daily reflections into infinity} blog post yet? If not, can you think of a better place for them? They look ugly under each blog post, in my current opinion (last confirmed 17 August 2026)

\item Blog: What's the biggest block to finishing one more of your in-progress blogs?

\item Touch type: Have you spent five minutes touch typing today? If not, do that now.

\item Physical Journal Writing: When did you last do this? How did the time feel? If it was more than 24 hours ago, why hasn't it happened again?

\item Physical Poem Writing: When did you last write a poem? How did the time feel? If more than 24 hours ago, why hasn't it happened again? What poem type was it, how did that feel and go? What aids and tools did you need?

\item Physical Writing: How's the new hand looking? Your current goals were that it is: printed, pretty (other goals), any progress on that front? Think using an n-graph map would help?

\item Rereading old blog posts: What's the most recent post you re-read? What's the new most recent post?

\item Rereading old journals: What journal are you on? What are you reading through now? Is it worth transfer to another medium? If so, where did you transfer it?

\item Reading other old writing: What writing are you reading? Anything worth maintaining? If not, why didn't you throw it away?

\item Reading old web serial: What chapter are you up to? Did you add notes to the remembrance document?

\item Regular embroidery practice: What's the current non-counted embroidery you're working on? When did you last work on it, how did the time feel? If more than 24 hours ago, why?  What do you think it would look like with five minutes of dedicated work? How close were you to the estimate?

\item Regular embroidery practice: What's the current counted embroidery you're working on? When did you last work on it, how did the time feel? If more than 24 hours ago, why? What do you think it would look like with five minutes of dedicated work? How close were you to the estimate?

\item Prayer habit: When's the last time you spent in prayer? How'd the timing feel? How'd the length feel? What did you pray? If more than 8 hours ago, why haven't you prayed since?

\item Jewish identity: What's one medium thing you can do to embrace this heritage? Is it really smaller things?

\item Jewish identity: What's at least one small thing you can do to embrace heritage? Do it.

\item What time are you ending this reflection? How's the time feel? How'd reflecting go?

\end{enumerate}

I also sometimes feel like the daily reflection is... far too long for me.  
What's the \say{I have no spoons but need to eat some cereal} version of the daily reflection?

Looking at the items (from the previous list:

\begin{itemize}

\item Absolutely keep 1  
\item Don't really need 2 or 3 or 4  
\item Probably keep 5, 6  
\item Def keep 7  
\item Drop 8, 9  
\item 10 through 12 can just be combined to \say{music?}  
\item 13 14 15 16 can all be dropped because RebelFit is not a no-spoons activity  
\item 17 can be changed to just \say{stretch}  
\item 18 should stay. Good way to get some efficacy in  
\item 19 through 22 should really move up to the top few because keep blogging together  
\item 23 maybe keep, good easy thing to do  
\item 24 keep  
\item 25 keep  
\item 26 drop  
\item 27 keep maybe  
\item 28, 29, 30 drop  
\item 31 or 32, combine to just \say{embroider?}  
\item 33 reduce to just \say{pray}?  
\item drop 34, 35  
\item Keep 36

\end{itemize}

so full list becomes (after adding in some food, water):

\begin{enumerate}

\item What time are you starting this reflection? How's this time feel?

\item Blogging: when did you last blog? How did the time feel? How long did you spend?

\item Blogging: Any changes needed to the template? If so, why haven't you made them?

\item Blogging: Have you written and posted your \say{blogs-im-writing} blog post yet? Why not?

\item Blogging: Have you written and posted your \say{do a single time tasks} blog post yet? Why not?

\item Blogging: Have you written and posted your \say{this is literally just daily reflections into infinity} blog post yet? If not, can you think of a better place for them? They look ugly under each blog post, in my current opinion (last confirmed 17 August 2026)

\item Blogging: What's the biggest block to finishing one more of your in-progress blogs?

\item Daily reflection: Any questions needed to this template? If so, what and why?

\item Home: What's one large thing that home setup needs right now? Is there a way to break it into small and medium tasks you can do?

\item Home: What's at least one medium thing you could do today to improve your home setup? Is there a way to break it or them into small tasks?

\item Home: What's at least one small thing you can do today to improve the shape of your home? If multiple, what is the one you're going to pause right now to do, and what is the most important of the remaining?

\item Career aspirations: Any changes?

\item If doing serials, look above for questions

\item Music: Have you sung today? If not, do a vocal warmup

\item Music: Have you played guitar today? If not, do at least one hand's flexibility exercise.

\item Music: What is the smallest possible piece of composition you could do right now that would feel like it \say{counts} to you? Do that now if you haven't already composed today.

\item RebelFit: What's in the way of it being finished? Can you break it into large, medium, small tasks?

\item RebelFit: What's at least one large task? Can you break it into medium and small tasks?

\item RebelFit What's at least one medium task? Can you break it or them into small tasks?

\item RebelFit: What's at least one small task? Can it be done within two minutes? If not, it's not small. If so, go do it.

\item Stretching: When were your last three stretches? What routine (lengths) were they? How did the time feel for stretching? How were the routines? If it's been more than 36 hours since three have been done, what's in the way?

\item Cleaning inbox: What's the inbox count? Go move or delete 25 emails, based on whether they need to be kept or not. New inbox count?

\item Touch type: Have you spent five minutes touch typing today? If not, do that now.

\item Physical Journal Writing: When did you last do this? How did the time feel? If it was more than 24 hours ago, why hasn't it happened again?

\item Physical Poem Writing: When did you last write a poem? How did the time feel? If more than 24 hours ago, why hasn't it happened again? What poem type was it, how did that feel and go? What aids and tools did you need?

\item Physical Writing: How's the new hand looking? Your current goals were that it is: printed, pretty (other goals), any progress on that front? Think using an n-graph map would help?

\item Rereading old blog posts: What's the most recent post you re-read? What's the new most recent post?

\item Rereading old journals: What journal are you on? What are you reading through now? Is it worth transfer to another medium? If so, where did you transfer it?

\item Reading other old writing: What writing are you reading? Anything worth maintaining? If not, why didn't you throw it away?

\item Reading old web serial: What chapter are you up to? Did you add notes to the remembrance document?

\item Regular embroidery practice: What's the current non-counted embroidery you're working on? When did you last work on it, how did the time feel? If more than 24 hours ago, why? What do you think it would look like with five minutes of dedicated work? How close were you to the estimate?

\item Regular embroidery practice: What's the current counted embroidery you're working on? When did you last work on it, how did the time feel? If more than 24 hours ago, why? What do you think it would look like with five minutes of dedicated work? How close were you to the estimate?

\item Prayer habit: When's the last time you spent in prayer? How'd the timing feel? How'd the length feel? What did you pray? If more than 8 hours ago, why haven't you prayed since?

\item Jewish identity: What's one medium thing you can do to embrace this heritage? Is it really smaller things?

\item Jewish identity: What's at least one small thing you can do to embrace heritage? Do it.

\item Eaten recently enough? Eaten enough recently? If no to either, eat.

\item Drunk recently enough? Drunk enough recently? If no to either, drink.

\item What time are you ending this reflection? How's the time feel? How'd reflecting go?

\end{enumerate}

and the short version:

\begin{enumerate}

\item What time are you starting this reflection? How’s this time feel?

\item OPTIONAL Home: What’s one large thing that home setup needs right now? Is there a way to break it into small and medium tasks you can do?

\item OPTIONAL Home: What’s at least one medium thing you could do today to improve
your home setup? Is there a way to break it or them into small tasks?

\item Home: What’s at least one small thing you can do today to improve the
shape of your home? If multiple, what is the one you’re going to pause
right now to do, and OPTIONAL what is the most important of the remaining?

\item Have you done music today? If not, can you sing to a song right now, even if just for a  moment?

\item Have you stretched today? If not, go touch your toes for a few breaths and then do a sun salutation.

\item Delete 25 emails from inbox. Can go to random spot for easy ones (or just search for something you know you can delete)

\item Touch type: Have you spent five minutes touch typing today? If not, do that now.

\item Physical journal: if you haven't done today, go do

\item Physical Poem Writing: if you haven't done today, go do

\item OPTIONAL Rereading old blog posts: What’s the most recent post you re-read? What’s the new most recent post?

\item If you haven't embroidered today, go do a few stitches

\item What's a prayer you feel right now? If nothing, go pray the St. Michael prayer.

\item What time are you ending this reflection? How’s the time feel? How’d reflecting go?

\end{enumerate}

Since I have spoons, I'll now add the reflection for today to the bottom of the list, and then hit post on this blog

\draft{1: 16 August 2026}

I've been realizing lately that I feel somewhat adrift.

That's a great and normal feeling; most often, this is a sign that I'm recovering from a deep well of depression.\footnote{initially I wrote sorrow, but that's not entirely accurate. Sorrow comes from disappointment and loss, not just from neurochemical breaking.}

However, it's also something that I still feel the need to address.  
Right now is my best time for improving the way that I am.  
So, what is the way that I want to change my life?

Or, more accurately, what are the \textit{ways} that I want to do so?

\footnote{wow I can never decide between itemize and enumerate. I think that in general life I use enumerate because nice to be able to uniquely identify with the number. In TeX (\LaTeX for the pedants), though, I do find that I tend towards non-enumeration because I don't think of what I'm doing as particularly ordered.  
The items are just a list of what I need, rather than any sort of sorted list.  
May be worth considering in the future}

\begin{itemize}

\item Figure out this blog.

What does that mean?  
\begin{itemize}  
\item Revamp\footnote{will that word appear too much in this post? To some people, the twice I've used it so far is already too much} the upcoming posts so that it's instead a reminder to myself of the posts that I've started but not yet finished.

I'll do that after this post is finished.

\item Revamp the daily reflection. I don't really think that the current set of questions serve me. That's something I'll plan to do after making this list

\item Figure out when in my day it makes the most sense to write, both for these next two weeks and going forward.  
\end{itemize}

\item Figure out a way to maintain a home that I can invite others into.

Honestly, that doesn't just require maintenance or even repair\footnote{not in the normal usage, to be clear. Repair here meaning like \say{to return the state}, not \say{fix a broken thing}}, but also finishing moving in in some regards.  
Let's think of some things that are needing to have a routine:

(N.B. I realized that right now I have my list of things that need a routine elsewhere, and right now what I moreso need is both to do the tasks and to get the apartment set up)

\item Get the career aspirations settled.

I've known that I don't want to stay at my current place forever. My partner moving away did little except exacerbate that urge to leave.\footnote{to be clear, did little to that knowledge. Her moving away did a lot to many things, most of which feel negative to me because I love being near her. They are positive towards the future we want to build, however}

What does this require?  
\begin{itemize}

\item Job search with more intention.  
\item Make my resume better, and probably make it easier to machine scan.\footnote{I do now understand the frustration others have of \say{I upload the resume and then I manually type the entire resume in again. The fact that the system my last application went to doesn't allow for deleting intermediate entries, only the most recent one, is incredibly frustrating}}

\end{itemize}

\item Decide how I'm feeling about the web serials right now.

In general, I want to be writing at least one of them, but I haven't.  
Why not?

I'm not totally sure.

That's something that I'll need to consider in the future, but my best guess is that i don't really want to do much, and it's a low enough priority item that it's fallen off of every day's list.  
Anyways, that's something else I'll need to figure out in the next few days.

\item Get RebelFit finished.

That's something I've been meaning to do for almost a full year now.\footnote{woot happy 1 year anniversary of becoming Dr.!}  
What's the barrier?

Primarily the barrier is that I want to revamp most of the code, now that I know a little bit more about objective programming.\footnote{or object-oriented. I'm sure that there's a difference and that I probably mean the orientation rather than the modifier}

\item Find a way to get stretching into my daily routine.

Pilates has been great for me and my body, but it's not everything, especially when I push myself\footnote{potentially too} hard.  
Stretches feel good, and that feels good.

\item Get music into my daily routine, both for the next few weeks and going forward.

I want to make music more than just the approximately weekly times that I meet with a friend.  
That can mean many things, but almost certainly requires me playing my guitar.  
Right now the guitar doesn't speak to me, which is almost always a sign of needing new strings, and so that's probably the first step.

 \end{itemize}

So, what else do I need to do to revamp my life?

That feels like as good of a starting place as any.

Honestly, this doesn't necessarily feel like as good of a post as it could be.  
Then again, if the blog is a diary\footnote{as I discussed in large part some long time ago}, then the more introspective aspect of this blog post is probably actually exactly what I want and need.  
Eh, it's something that I can post, if nothing else.

\dailyreflect{17 August 2026}

\begin{enumerate}

\item What time are you starting this reflection? How's this time feel?

13:42. A little late, given that I've been blogging for a few hours. I also feel hot and sweaty in current location, so let's use that as the motivation I need to go stand up and move to a different space.

Much better

\item Blogging: when did you last blog? How did the time feel? How long did you spend?

Todayish, started approx 11:30, I think? Wild, two hours goes by so fast. Time felt nice, but I also got up later than usual today because no Pilates.  
Long term, a two hour block of writing time is not super compatible with the rest of life as I will again live.

\item Blogging: Any changes needed to the template? If so, why haven't you made them?

Not any longer for now?

\item Blogging: Have you written and posted your \say{blogs-im-writing} blog post yet? Why not?

Nope! Today's post is this one, and tomorrow's is hopefully that.

\item Blogging: Have you written and posted your \say{do a single time tasks} blog post yet? Why not?

Nope! Today's is this, overmorrow's is hopefully that.

\item Blogging: Have you written and posted your \say{this is literally just daily reflections into infinity} blog post yet? If not, can you think of a better place for them? They look ugly under each blog post, in my current opinion (last confirmed 17 August 2026)

Nope! TBH, I don't know how ugly it really is.

\item Blogging: What's the biggest block to finishing one more of your in-progress blogs?

Uhhhhh at this point I don't really know what all of them are.   
Otherwise, block is that I need to be off writing time for a little now.

\item Daily reflection: Any questions needed to this template? If so, what and why?

Nope!  
Though, this q should go above the other blogging questions.\footnote{and so I make that change now}

\item Home: What's one large thing that home setup needs right now? Is there a way to break it into small and medium tasks you can do?

Right now I think that the big thing would be clearing the music room floor.  
That can be broken into medium task:

\begin{itemize}

\item Clear a path through music room floor

\item Throw out trash on floor

\item Move things into the other rooms as needed to make space for now.

\end{itemize}

\item Home: What's at least one medium thing you could do today to improve your home setup? Is there a way to break it or them into small tasks?

From above: Clear a path through music room floor.  
Only real small way that I can make that happen is to just list each item individually

\item Home: What's at least one small thing you can do today to improve the shape of your home? If multiple, what is the one you're going to pause right now to do, and what is the most important of the remaining?

Clear a space to walk into music room.  
Not going to do that now, because I'm planning to do the medium task when I finish this post.

\item Career aspirations: Any changes?

Not really.

\item If doing serials, look above for questions

Whoop! For now, let's change to \say{serials? Thoughts?}  
Hmmm, that's probably something that implies I want the blogging answers at the bottom.

Let's do that in the template as well.

\item Music: Have you sung today? If not, do a vocal warmup

Nope! Vocal warmup will be done in car maybe.

\item Music: Have you played guitar today? If not, do at least one hand's flexibility exercise.

Hands are sweaty right now and I don't want to mess with the guitar since it has new strings.

\item Music: What is the smallest possible piece of composition you could do right now that would feel like it \say{counts} to you? Do that now if you haven't already composed today.

Honestly, writing down a melody in MuseScore would hit that urge for me. That's not something I'm going to do right now, though.

\item RebelFit: What's in the way of it being finished? Can you break it into large, medium, small tasks?

Uhhhh pass.

\item RebelFit: What's at least one large task? Can you break it into medium and small tasks?

Rewrite the code. Smaller large tasks are rewrite the bounding algorithm and rewrite the point selection algorithm.  
Medium sized ones for point selection are each constant in the \say{pure} form, which is what will be sampled.

\item RebelFit What's at least one medium task? Can you break it or them into small tasks?

Any given \say{pure} constant, for now let's say picking C given I.

Also let's see if I can figure out how to pick sigma, C, I as my three base variables again.  
Or, at least, use A, B, C, I, sigma, figure out how to constrain using all of them.

One small task is finding the equation for C (normal) given A reduced constants and then using the bounds that I'll have had in A reduced form, making that.

\item RebelFit: What's at least one small task? Can it be done within two minutes? If not, it's not small. If so, go do it.

See above. Not doing now, though.

\item Stretching: When were your last three stretches? What routine (lengths) were they? How did the time feel for stretching? How were the routines? If it's been more than 36 hours since three have been done, what's in the way?

... no comment.

I found my packet form back when I was a diver, some of those might be good to do agin.

\item Cleaning inbox: What's the inbox count? Go move or delete 25 emails, based on whether they need to be kept or not. New inbox count?

Inbox count: 152003. 25 gets me basically to 151975, so let's go. Woot, that went fast enough. In the future, may need to remind myself that like, scrolling down through posts I'm not sure if I should delete still good.

\item Touch type: Have you spent five minutes touch typing today? If not, do that now.

Nope! Will maybe do after, after I get some non-screen tiem.

\item Physical Journal Writing: When did you last do this? How did the time feel? If it was more than 24 hours ago, why hasn't it happened again?

.... Anyways

\item Physical Poem Writing: When did you last write a poem? How did the time feel? If more than 24 hours ago, why hasn't it happened again? What poem type was it, how did that feel and go? What aids and tools did you need?

... Anyways.

\item Physical Writing: How's the new hand looking? Your current goals were that it is: printed, pretty (other goals), any progress on that front? Think using an n-graph map would help?

... Anyways

\item Rereading old blog posts: What's the most recent post you re-read? What's the new most recent post?

...Anyways

\item Rereading old journals: What journal are you on? What are you reading through now? Is it worth transfer to another medium? If so, where did you transfer it?

...Anyways

\item Reading other old writing: What writing are you reading? Anything worth maintaining? If not, why didn't you throw it away?

...Anyways

\item Reading old web serial: What chapter are you up to? Did you add notes to the remembrance document?

...Anyways

\item Regular embroidery practice: What's the current non-counted embroidery you're working on? When did you last work on it, how did the time feel? If more than 24 hours ago, why? What do you think it would look like with five minutes of dedicated work? How close were you to the estimate?

...Anyways

\item Regular embroidery practice: What's the current counted embroidery you're working on? When did you last work on it, how did the time feel? If more than 24 hours ago, why? What do you think it would look like with five minutes of dedicated work? How close were you to the estimate?

...Anyways

\item Prayer habit: When's the last time you spent in prayer? How'd the timing feel? How'd the length feel? What did you pray? If more than 8 hours ago, why haven't you prayed since?

...Anyways... Just did St. Michael in a prayerful attitude.

\item Jewish identity: What's one medium thing you can do to embrace this heritage? Is it really smaller things?

Uhhhhh. pass

\item Jewish identity: What's at least one small thing you can do to embrace heritage? Do it.

Pass

\item Eaten recently enough? Eaten enough recently? If no to either, eat.

Generally yes, but probably good to have an apple

\item Drunk recently enough? Drunk enough recently? If no to either, drink.

Generally yes, but given the sweat, take a drink.

\item What time are you ending this reflection? How's the time feel? How'd reflecting go?

13:58, so 16 minutes. Time feels generally good. Don't love how much reflecting is just going \say{pass}, today, but that's a downside of starting so many words into a document

\end{enumerate}

\end{document}